%%%%%%%%%%%%
%
% $Autor: Siddhanth Sharma $
% $Datum: 2026-01-02 $
% $Pfad: report/Contents/General/Display.tex $
% $Version: 4250 $
% !TeX spellcheck = en_GB/de_DE
% !TeX encoding = utf8
% !TeX root = ../../System/EdgeComputer/jetsonNano.tex
% !TeX TXS-program:bibliography = txs:///biber
%
%%%%%%%%%%%%


\chapter{Display as System Actor}
\label{chap:display_actor}

\section{Introduction}
In embedded \ac{AI} systems, \textbf{actors} are components that translate internal computation into observable system responses.
In this project, the actor is a \textbf{visual user interface} that presents the results of the license-plate pipeline on an \ac{HDMI}-connected monitor.
Instead of driving a physical actuator (gate, buzzer), the Jetson Nano runs a local dashboard that shows (i) an annotated live camera stream and (ii) the current session/billing state derived from the persistent database.

The display is connected to the Jetson Nano via the \textbf{\ac{HDMI} 2.0 (J3)} port on the Seeed Studio J101 carrier board.
It serves as a direct, on-device output interface for the operator during runtime (and is also useful during development).

\begin{figure}[H]
	\centering
	\includegraphics[width=0.95\linewidth]{Display/Display.png}
	\caption{Display actor: the Jetson Nano renders the local dashboard on an \ac{HDMI} monitor, combining live visual feedback and operational controls.}
	\label{fig:display_overview}
\end{figure}

\section{General}
Displays are standard output actors in access-control and surveillance systems because they enable immediate, human-readable verification of system behaviour.
In this setup, the \ac{HDMI}-connected monitor shows a \textbf{local web dashboard} running on the Jetson Nano.
The dashboard is opened in a browser \textit{on the Jetson itself} and continuously updates while the inference pipeline runs.

The operator sees the processed camera image (license plate region highlighted) and a structured list of recent detections and their derived state (timestamps and fee).
This keeps all output local to the edge device.

\section{Specific Actor}
\label{sec:display_specs}
In this context, a standard \ac{HDMI} display connected through the Jetson Nano’s \ac{HDMI} 2.0 Type-A port acts as the system actor.
The monitor renders a browser-based dashboard (served locally by the Jetson) that presents:
\begin{itemize}
	\item a live annotated video stream (frame overlays drawn by the inference pipeline), and
	\item a database-backed parking/session table for the most recent recognized plates.
\end{itemize}

\begin{table}[H]
	\centering
	\caption{Display module specifications and connection parameters \cite{seeedJ101Manual, jetson_nano_doc}}
	\label{tab:display_specs}
	\begin{tabular}{|p{5cm}|p{8cm}|}
		\hline
		\textbf{Parameter} & \textbf{Specification} \\ \hline
		Display Interface & \ac{HDMI} 2.0 Type-A (J3 port on carrier board) \\ \hline
		Supported Resolution & Up to 3840 × 2160 @ 60 Hz (4K) \\ \hline
		Color Depth & 24-bit RGB \\ \hline
		Power Supply & External 5 V DC adapter (display dependent) \\ \hline
		Connection Type & \ac{HDMI} male-to-male cable \\ \hline
		Display Technology & LCD/LED compatible \\ \hline
		Usage in Project & Local dashboard UI: live stream, session log, and operational actions \\ \hline
	\end{tabular}
\end{table}

\section{Specification}
The display acts purely as an output peripheral and requires no electrical modification.  
The only connection is the \ac{HDMI} cable carrying both video and synchronization signals from the Jetson Nano.  
Circuit representation is shown later in the overall system hardware diagram.

The complete system-level data flow between the camera, Jetson Nano, and display actor is illustrated in Figure~\ref{fig:dataflow_display}.

\begin{figure}[H]
	\centering
		\begin{tikzpicture}[
		node distance=2.2cm,
		>=latex,
		thick,
		align=center,
		font=\sffamily
		]
		
		% Nodes (top to bottom)
		\node[draw, fill=blue!10, rounded corners,
		minimum width=6cm, minimum height=1.2cm]
		(camera) {\textbf{Microsoft LifeCam Show}\\(Sensor Input)};
		
		\node[draw, fill=green!10, rounded corners,
		minimum width=6cm, minimum height=1.2cm,
		below=of camera]
		(jetson) {\textbf{NVIDIA Jetson Nano}\\\textbf{(Seeed Studio reComputer J1010)}\\(AI Processing Unit)};
		
		\node[draw, fill=orange!10, rounded corners,
		minimum width=6cm, minimum height=1.2cm,
		below=of jetson]
		(display) {\textbf{HDMI Display}\\(System Actor Output)};
		
		% Arrows
		\draw[->, thick] (camera) -- node[right, align=center]{\textit{Video Stream}\\(USB 2.0)} (jetson);
		\draw[->, thick] (jetson) -- node[right, align=center]{\textit{Processed Output}\\(HDMI Signal)} (display);
		
		% Optional feedback arrow
		\draw[<-, dashed, gray, thick]
		(camera.west) .. controls +(-3,-1) and +(-3,1) ..
		node[left, align=center, text=gray!70, font=\itshape]
		{Control Signal\\/ Status Log} (jetson.west);
		
	\end{tikzpicture}
	\caption{Vertical data flow of the system: image acquisition from the LifeCam,\ac{AI} inference on Jetson Nano, and output display via \ac{HDMI}.}
	\label{fig:dataflow_display}
\end{figure}

\section{Library}
\subsection{Description}
From a hardware perspective, the display actor relies on the Jetson Nano graphics output stack to generate a stable \ac{HDMI} signal.
The operating system detects the connected monitor through \textbf{EDID} (Extended Display Identification Data), negotiates a compatible resolution/refresh rate, and continuously drives the \ac{HDMI} link during runtime.
The display therefore requires no dedicated device driver integration beyond the standard Jetson/Linux graphics stack; it behaves as an external \ac{HDMI} sink.

\subsection{Installation}
No special installation is required for the display hardware itself beyond connecting the \ac{HDMI} cable and ensuring the Jetson graphics stack is active.
For verification and diagnostics, the monitor connection can be inspected from the shell:

\begin{tcolorbox}[colback=white,colframe=black,coltext= ShellColor, sharp corners, boxrule=0.8pt]
	\begin{small}
		\begin{verbatim}
			
			xrandr	\end{verbatim}
	\end{small}
\end{tcolorbox}

\subsection{Functions}
Essential software elements used for the display actor:
\begin{itemize}
	\item \textbf{EDID-based detection:} the connected monitor is recognized and a display mode is selected automatically.
	\item \textbf{Mode control and diagnostics:} resolution and refresh rate can be inspected and adjusted via \SHELL{xrandr}.
	\item \textbf{Continuous video output:} the Jetson renders the user interface to the framebuffer, which is output over \ac{HDMI}.
\end{itemize}

\subsection{Data Characteristics}
\paragraph{Data Types:} RGB camera frames (\texttt{uint8} arrays) that are encoded as JPEG bytes for streaming; JSON objects for table updates; SQLite rows for persistent session state.

\paragraph{Data Structure:} The dashboard consumes (i) a continuous MJPEG byte stream for the live view and (ii) periodic JSON responses containing plate strings, timestamps, and computed fee values.

\paragraph{Data Quantity:} The UI updates continuously for the video stream and polls the JSON endpoint periodically (e.g., every second) to refresh the table.

\paragraph{Data Quality:} The video overlays provide immediate visual validation of plate localization, while the table view provides readable, persistent state for the operator.
The \ac{HDMI} output ensures stable on-device presentation without requiring network connectivity.

\section{Example}
In this section, a complete example illustrates how the display actor visualizes the system response.

\subsection{Manual}
To operate the display actor, connect the \ac{HDMI} cable between the Jetson Nano and monitor, select the proper input source, and power both devices.
Start the Jetson application and ensure the user interface is rendered on the connected monitor.
For demonstration/kiosk operation, use full-screen mode and ensure that the monitor resolution is set so all UI elements are visible.

\subsection{Inside the Example and Code}
The rendered interface on the monitor is organized as a two-panel layout:
\begin{itemize}
	\item \textbf{Left:} live camera stream with overlays (bounding box on the detected license plate).
	\item \textbf{Right:} \textit{PARKING LOG (DB)} table showing the most recent recognized plate(s), timestamps, and the computed fee derived from persisted visit data.
\end{itemize}

\begin{figure}[H]
	\centering
	\includegraphics[width=0.95\linewidth]{Display/dashboard.png}
	\caption{Full dashboard on the \ac{HDMI} display: annotated live stream (left) and database-backed log table (right).}
	\label{fig:display_dashboard}
\end{figure}

\begin{figure}[H]
	\centering
	\begin{subfigure}[t]{0.48\linewidth}
		\centering
		\includegraphics[width=\linewidth]{Display/sidebar.png}
		\caption{Table panel (\textit{PARKING LOG (DB)}) with fee output.}
		\label{fig:display_sidebar}
	\end{subfigure}
	\hfill
	\begin{subfigure}[t]{0.48\linewidth}
		\centering
		
\includegraphics[width=\linewidth]{Display/shutDownbutton.png}
		\caption{Shutdown control for a controlled poweroff from the UI.}
		\label{fig:display_shutdown}
	\end{subfigure}
	\caption{Key dashboard UI elements used as the system actor.}
	\label{fig:display_ui_elements}
\end{figure}


\subsection{Code}
No dedicated code is required for the display hardware itself beyond the standard \ac{HDMI} output of the Jetson Nano.
Software implementation details of the dashboard/UI are documented in the application deployment chapters; here the focus remains on the physical actor and its integration.

\section{Calibration}
Calibration is limited to brightness, contrast, and resolution adjustments for optimal visibility.  
Settings are managed using \texttt{xrandr} or the desktop display configuration of the Jetson runtime environment.
For kiosk-like operation, browser zoom and full-screen mode can be used to match the monitor resolution.

\section{Simple Code}
The minimal functional verification for the display actor is to confirm that the monitor is detected and a valid mode is active (resolution/refresh rate), and that the Jetson renders output continuously to the screen.
This is typically verified using \SHELL{xrandr} together with a visual check that the UI is visible without flicker or tearing.

\section{Simple Application}
During live inference, the display actor provides:
\begin{itemize}
	\item \textbf{Live verification:} annotated camera frames with the detected plate region highlighted.
	\item \textbf{Operational table:} a continuously updated log of plates and their derived state from the SQLite database.
	\item \textbf{Operator control:} a shutdown action for safe poweroff after the demonstration.
\end{itemize}
This provides clear user feedback and makes system behaviour observable directly at the Jetson device.

\section{Tests}
\subsection{Simple Function Test}
Run:

\begin{tcolorbox}[colback=white,colframe=black,coltext= ShellColor, sharp corners, boxrule=0.8pt]
	\begin{small}
		\begin{verbatim}
			
			xrandr	\end{verbatim}
	\end{small}
\end{tcolorbox}
to confirm that the \ac{HDMI} display is detected and configured.

\subsection{Test all Functions}
Execute the complete dashboard application and verify:
\begin{itemize}
	\item the dashboard loads in the on-device browser,
	\item the live stream is visible and updates continuously, and
	\item the table updates from SQLite and the shutdown button triggers the expected action (only when confirmed by the operator).
\end{itemize}

\section{Constraints}
\begin{itemize}
	\item Requires \ac{HDMI}-compatible display; analog VGA unsupported.  
	\item Browser must run on the Jetson Nano to display the local dashboard.  
	\item Video overlays must avoid covering key regions needed for visual validation.  
	\item Manual brightness/contrast tuning recommended.  
\end{itemize}

\section{Further Readings}
\begin{itemize}
	\item NVIDIA Jetson Nano Developer Guide — \cite{jetson_nano_doc}
	\item Seeed Studio reComputer J101 Manual — \cite{seeedJ101Manual}
	\item OpenCV Video Display Documentation — \cite{opencv_doc_2024}
	\item Embedded Edge \ac{AI} Visualization Techniques — \cite{edgeAIbook}
\end{itemize}
